\chapter{Análisis del Módulo TLR}
\label{cap:analisis_apollo}

Este capítulo describe el proceso de análisis del módulo de detección de semáforos (\textit{Traffic Light Recognition}, TLR) de Apollo. El objetivo del análisis no fue documentar el módulo por sí mismo, sino reunir el conocimiento necesario para poder aislarlo: entender qué hace, cuáles son sus límites, qué necesita del resto del sistema y qué simplificaciones serían aceptables.

El capítulo sigue ese orden. Primero plantea las preguntas que guiaron el análisis. Después responde a alto nivel cómo está organizado el módulo dentro de Apollo y qué hacen sus etapas, acompañado de un diagrama de bloques y un diagrama de secuencia con \textit{input}/\textit{output}. Luego presenta los hallazgos del código fuente que resultaron determinantes para el aislamiento. Y cierra con las decisiones de aislamiento que se desprenden de todo lo anterior. El detalle técnico de cómo funciona cada etapa internamente se posterga al Capítulo~\ref{cap:funcionamiento_tlr}, donde tiene espacio propio.

\section{Preguntas que Guiaron el Análisis}
\label{sec:preguntas_analisis}

El análisis estuvo orientado por cuatro preguntas. Cada una apunta a una decisión concreta de aislamiento.

\begin{enumerate}
    \item \textbf{¿Cuál es el alcance funcional del módulo TLR?} \\
    ¿Qué hace exactamente? ¿Cuáles son sus entradas y salidas? ¿De qué otros módulos depende para funcionar?

    \item \textbf{¿Qué partes del código son esenciales y cuáles son accesorias?} \\
    ¿Qué es lógica \textit{core} del TLR (decisiones, parámetros, modelos) y qué es infraestructura del entorno Apollo (CyberRT, Bazel, integración con otros módulos)? El aislamiento depende de poder hacer esta distinción.

    \item \textbf{¿Qué parámetros y reglas de decisión utiliza el sistema?} \\
    ¿Cuáles son los \textit{thresholds}, pesos y reglas de negocio que determinan el comportamiento observable? Estos son los que el módulo aislado debe reproducir.

    \item \textbf{¿Qué simplificaciones podemos hacer sin afectar la lógica \textit{core}?} \\
    ¿Qué dependencias podemos eliminar manteniendo equivalencia funcional? ¿Qué se puede reemplazar por una interfaz más simple sin perder la semántica del módulo?
\end{enumerate}

Las dos primeras preguntas delimitan qué entra al módulo aislado y qué queda afuera. La tercera define qué se debe preservar. La cuarta habilita la modularización: identifica qué piezas del entorno original pueden ser reemplazadas por equivalentes más simples sin que el comportamiento del módulo cambie.

\section{El Módulo TLR dentro de Apollo}

\subsection{Ubicación en el Código}

El TLR no es un único componente monolítico, sino un conjunto de cuatro subcomponentes que viven en \texttt{modules/perception/} del repositorio de Apollo:

\begin{verbatim}
modules/perception/
|-- traffic_light_region_proposal/  # Generacion de ROIs desde HD-Map
|-- traffic_light_detection/        # Deteccion de semaforos en imagen
|-- traffic_light_recognition/      # Clasificacion de color
`-- traffic_light_tracking/         # Tracking temporal y reglas
\end{verbatim}

Cada subcomponente es un \textit{component} en el sentido de Apollo: tiene sus propios archivos de configuración (\texttt{.pb.txt}, \texttt{.proto}), sus reglas de \textit{build} para Bazel (\texttt{BUILD}), su descripción de despliegue (\texttt{cyberfile.xml}, \texttt{.dag}, \texttt{.launch}) y su código fuente. Los cuatro componentes se ejecutan en cadena dentro del proceso de percepción del vehículo, comunicándose a través del \textit{framework} de mensajería CyberRT.

\subsection{Infraestructura de la que Depende}

Para que los cuatro componentes funcionen, Apollo necesita además infraestructura que es transversal al sistema completo:

\begin{itemize}
    \item \textbf{CyberRT:} \textit{framework} de mensajería propio que orquesta la comunicación entre componentes a través de \textit{channels}. Sin CyberRT no hay forma estándar de pasar un \textit{frame} de la cámara al componente de detección.
    \item \textbf{Map module:} acceso al HD-Map. El componente de \textit{region proposal} consulta el HD-Map en tiempo real para saber dónde hay semáforos cerca del vehículo.
    \item \textbf{Perception common:} utilidades compartidas con el resto del módulo de percepción: \textit{inference engines}, transformaciones de coordenadas, primitivas geométricas.
    \item \textbf{Modules common:} utilidades transversales del sistema (\textit{math}, \textit{configs}, \textit{logging}).
    \item \textbf{Calibración de cámara:} parámetros intrínsecos y extrínsecos del vehículo que se usan para proyectar coordenadas 3D del mundo a píxeles de imagen.
    \item \textbf{Caffe (vía libtorch):} \textit{framework} de \textit{deep learning} sobre el que corren los cuatro modelos del TLR. Requiere CUDA y una GPU NVIDIA específica.
\end{itemize}

El TLR no puede ejecutarse sin esta infraestructura; cualquier estrategia de aislamiento tiene que decidir qué se hace con ella.

\section{Cómo Funciona el Módulo: Vista a Alto Nivel}

\subsection{Las Cinco Etapas}

El pipeline del TLR procesa cada \textit{frame} de la cámara en cinco etapas secuenciales:

\begin{enumerate}
    \item \textbf{Preprocesamiento (\textit{Region Proposal}).} A partir de la posición del vehículo y el HD-Map, calcula qué semáforos físicos deberían aparecer en la imagen y proyecta sus coordenadas 3D del mundo a regiones rectangulares 2D en la imagen (\textit{projection boxes}). También decide qué cámara usar cuando hay varias disponibles. Cada semáforo identificado mantiene un \textit{signal\_id} persistente proveniente del HD-Map.

    \item \textbf{Detección.} Para cada \textit{projection box}, recorta una región expandida de la imagen y la pasa por una red convolucional (estilo \textit{Faster R-CNN}) que devuelve detecciones de semáforos con un \textit{bounding box} y un tipo (vertical, horizontal, \textit{quad}). Las detecciones no llevan identidad: son simplemente recuadros en la imagen.

    \item \textbf{Asignación.} Vincula cada detección de la red con la \textit{projection box} que mejor le corresponde. Usa el algoritmo húngaro sobre una matriz de costos que combina proximidad espacial y confianza de detección. Su salida son pares ($i$, $j$) que dicen ``la \textit{projection box} $i$ corresponde a la detección $j$'', y por lo tanto a un \textit{signal\_id} concreto.

    \item \textbf{Reconocimiento.} Para cada semáforo detectado, recorta el \textit{bounding box} de la detección y lo pasa por una de tres redes convolucionales especializadas (una por orientación: vertical, horizontal, \textit{quad}). La red devuelve probabilidades para cuatro estados: apagado (\textit{black}), rojo, amarillo y verde.

    \item \textbf{Tracking.} Mantiene un historial por \textit{signal\_id} a través de los \textit{frames}. Aplica reglas de consistencia temporal y de seguridad: transiciones de color inválidas se corrigen, detecciones aisladas de ``apagado'' no destruyen el color anterior, y se detecta parpadeo en verde con una ventana temporal específica.
\end{enumerate}

La Figura~\ref{fig:pipeline_alto_nivel} muestra las cinco etapas como bloques. La Figura~\ref{fig:diagrama_secuencia} las muestra como un diagrama de secuencia donde cada flecha indica el \textit{input}/\textit{output} que pasa de una etapa a la siguiente.

\subsection{Diagrama de Bloques}

\begin{figure}[H]
\centering
\begin{tikzpicture}[
    node distance=0.55cm,
    block/.style={rectangle, draw, rounded corners, minimum width=8.5cm, minimum height=0.9cm, align=center, font=\small},
    io/.style={rectangle, draw, dashed, rounded corners, minimum width=8.5cm, minimum height=0.8cm, align=center, font=\small\itshape},
    arrow/.style={-{Latex[length=2.5mm]}, thick}
]
    \node[io] (in) {Frame de cámara + Pose del vehículo + HD-Map};
    \node[block, below=of in] (preproc) {1. Preprocesamiento (\textit{Region Proposal})};
    \node[block, below=of preproc] (det) {2. Detección};
    \node[block, below=of det] (asg) {3. Asignación (algoritmo húngaro)};
    \node[block, below=of asg] (rec) {4. Reconocimiento};
    \node[block, below=of rec] (trk) {5. Tracking};
    \node[io, below=of trk] (out) {Estados de semáforos con \textit{signal\_id} persistente};

    \draw[arrow] (in) -- (preproc);
    \draw[arrow] (preproc) -- (det);
    \draw[arrow] (det) -- (asg);
    \draw[arrow] (asg) -- (rec);
    \draw[arrow] (rec) -- (trk);
    \draw[arrow] (trk) -- (out);
\end{tikzpicture}
\caption[Pipeline del módulo TLR de Apollo (5 etapas).]{Pipeline del módulo TLR de Apollo: cinco etapas secuenciales que transforman un \textit{frame} de cámara más información geográfica en un conjunto de estados de semáforos identificados.}
\label{fig:pipeline_alto_nivel}
\end{figure}

\subsection{Diagrama de Secuencia (Input/Output entre Etapas)}

\begin{figure}[H]
\centering
\resizebox{\textwidth}{!}{%
\begin{tikzpicture}[
    font=\small,
    actor/.style={rectangle, draw, fill=gray!10, text width=2.1cm, minimum height=0.7cm, align=center, inner sep=3pt},
    arrow/.style={-{Latex[length=2mm]}, thick},
    msg/.style={font=\scriptsize, align=center}
]
    \node[actor] (env) at (0,0) {\scriptsize Cámara\\\scriptsize + HD-Map};
    \node[actor] (p1) at (3,0) {Preprocesa-\\miento};
    \node[actor] (p2) at (6,0) {Detección};
    \node[actor] (p3) at (9,0) {Asignación};
    \node[actor] (p4) at (12,0) {Reconoci-\\miento};
    \node[actor] (p5) at (15,0) {Tracking};

    \foreach \x in {0,3,6,9,12,15} {
        \draw[dashed] (\x,-0.5) -- (\x,-9.5);
    }

    \node[msg, anchor=west] at (0.1,-1.0) {\textit{frame}, pose, \textit{signals} del HD-Map};
    \draw[arrow] (0,-1.3) -- (3,-1.3);

    \node[msg, anchor=west] at (3.1,-2.1) {M \textit{projection boxes} con \textit{signal\_id}};
    \draw[arrow] (3,-2.4) -- (6,-2.4);

    \node[msg, anchor=west] at (6.1,-3.2) {N detecciones (sin identidad)};
    \draw[arrow] (6,-3.5) -- (9,-3.5);

    \node[msg, anchor=west] at (3.1,-4.5) {(las M \textit{projection boxes} también)};
    \draw[arrow] (3,-4.8) -- (8.95,-4.8);

    \node[msg, anchor=west] at (9.1,-5.7) {K pares (\textit{proj\_idx}, \textit{det\_idx})};
    \draw[arrow] (9,-6.0) -- (12,-6.0);

    \node[msg, anchor=west] at (12.1,-7.0) {color por detección asignada};
    \draw[arrow] (12,-7.3) -- (15,-7.3);

    \node[msg, anchor=west] at (15.1,-8.2) {\textit{signal\_id} $\to$ (color revisado, \textit{blink})};
    \draw[arrow] (15,-8.5) -- (16.8,-8.5);
\end{tikzpicture}%
}
\caption[Diagrama de secuencia del pipeline TLR.]{Diagrama de secuencia del pipeline TLR: cada flecha representa el dato que cruza de una etapa a la siguiente. Las cantidades $M$ (cuántos semáforos hay en el HD-Map a esa distancia), $N$ (cuántas detecciones produjo la red) y $K \le \min(M,N)$ (cuántas asignaciones encontró el húngaro) varían por \textit{frame}.}
\label{fig:diagrama_secuencia}
\end{figure}

\subsection{Inputs y Outputs del Módulo Completo}

Visto como caja, el módulo recibe:

\begin{itemize}
    \item Un \textit{frame} de cámara (uno por instante, en tiempo real).
    \item La pose del vehículo en coordenadas del mundo.
    \item Acceso al HD-Map para hacer consultas dentro de un radio dado (típicamente $\sim$150\,m).
    \item Parámetros de calibración de las cámaras disponibles.
\end{itemize}

Y produce:

\begin{itemize}
    \item Un conjunto de semáforos identificados con \textit{signal\_id} persistente.
    \item Para cada uno: el color reconocido (\textit{black}/red/yellow/green) y un \textit{flag} de \textit{blink} cuando aplica.
\end{itemize}

El \textit{signal\_id} persistente es lo que permite que la etapa de \textit{tracking} mantenga un historial coherente \textit{frame} a \textit{frame}: el mismo semáforo físico siempre se identifica con el mismo \textit{id}, aunque cambien sus píxeles en la imagen.

\section{Hallazgos del Código Fuente}
\label{sec:hallazgos}

El análisis estático del código fuente reveló cuatro hallazgos que resultaron determinantes para definir el alcance del aislamiento.

\subsection{Hallazgo 1: Dos Tipos de Identificadores}

Apollo distingue de manera explícita dos identificadores que conviven en el pipeline:

\begin{table}[H]
\centering
\begin{tabular}{|l|p{5cm}|p{5cm}|}
\hline
\textbf{Tipo} & \textbf{Origen} & \textbf{Persistencia} \\
\hline
\textit{Signal\_id} & HD-Map (coordenadas 3D del mundo) & Permanente (mismo semáforo físico) \\
\hline
\textit{Detection\_id} & Asignado durante la detección en cada \textit{frame} & Solo válido para el \textit{frame} actual \\
\hline
\end{tabular}
\caption{Diferencias entre los dos tipos de identificadores que maneja el TLR.}
\end{table}

\textbf{Implicación:} el \textit{tracking} se apoya en el \textit{signal\_id}, no en la posición en píxeles. Esto significa que mover un semáforo en la imagen (por ejemplo entre \textit{frames}, por movimiento del vehículo) no confunde al sistema: la identidad viene del HD-Map, no del \textit{frame}. El módulo aislado tiene que preservar esa lógica: cualquiera sea la fuente de las \textit{projection boxes}, debe entregar \textit{signal\_id}s persistentes para que la etapa de \textit{tracking} funcione correctamente.

\subsection{Hallazgo 2: \textit{Semantic\_id} Diseñado pero No Implementado}

El análisis encontró código diseñado para una funcionalidad que en la práctica no se ejecuta:

\begin{codesnippet}
\begin{verbatim}
// traffic_light/region_proposal_component.cc:335
int cur_semantic = 0;  // Siempre 0, hardcodeado
for (auto& light : lights) {
    light->semantic = cur_semantic;
}
\end{verbatim}
\end{codesnippet}

El \textit{semantic\_id} debería agrupar semáforos de la misma intersección para aplicar \textit{voting} (si dos de tres semáforos del cruce ven verde, el tercero también debería verlo). El código de \textit{voting} existe en la etapa de \textit{tracking}, pero como \textit{semantic\_id} está hardcodeado a 0 para todos los semáforos, cada uno se trata como un grupo individual y el \textit{voting} se reduce a una operación trivial sobre un solo elemento.

\textbf{Implicación:} todos los semáforos se procesan independientemente, sin corrección por consistencia entre semáforos del mismo cruce. Esto es importante para el aislamiento por dos razones: el módulo aislado no necesita implementar el \textit{voting} efectivo (no aporta nada en el original), y este hueco constituye una oportunidad concreta de testing (qué pasa cuando un semáforo del mismo cruce contradice a otro).

\subsection{Hallazgo 3: Dependencia Crítica del HD-Map}

El HD-Map no solo provee las regiones donde buscar semáforos: es la fuente de identidad del sistema. Sin HD-Map:

\begin{itemize}
    \item No hay \textit{signal\_id}, por lo tanto no hay \textit{tracking} consistente.
    \item No hay \textit{projection boxes}, por lo tanto la red de detección no sabe dónde mirar.
    \item No hay reglas de selección de cámara, porque esas reglas se basan en la geometría 3D de los semáforos.
\end{itemize}

\textbf{Implicación para el aislamiento:} no podemos reproducir un HD-Map dinámico ni la proyección 3D→2D con calibración de cámara. Pero podemos \textit{simular su función} entregando \textit{projection boxes} pre-calculadas y \textit{signal\_id}s estables desde un archivo. Desde el punto de vista del resto del módulo, el \textit{input} es el mismo (una lista de \textit{bounding boxes} 2D con identidad persistente); la diferencia está en cómo se obtiene esa lista.

\subsection{Hallazgo 4: Selección Automática de Cámara}

Apollo opera con un sistema multi-cámara: una telephoto (25\,mm, $\sim$30° de campo de visión, mayor resolución a distancia) y una \textit{wide-angle} (6\,mm, $\sim$120° de campo de visión, mejor para semáforos cercanos). La etapa de preprocesamiento incluye lógica para elegir cuál usar en cada \textit{frame}: por defecto telephoto, pero \textit{wide-angle} si alguna \textit{projection box} queda fuera del campo de la telephoto.

\textbf{Implicación:} esta lógica depende del HD-Map (necesita las proyecciones para decidir) y del parque de cámaras del vehículo. Para el aislamiento, asumir una sola cámara no afecta la lógica de las etapas posteriores: ellas reciben una imagen y una lista de \textit{projection boxes}, sin importar de cuál cámara vinieron.

\section{Decisiones de Aislamiento}

Los hallazgos anteriores informan tres decisiones que definen qué se aísla y qué no.

\subsection{Alcance Funcional}

\textbf{Incluir} (la lógica \textit{core} del TLR):
\begin{itemize}
    \item Las cinco etapas del pipeline.
    \item Los cuatro modelos de red neuronal (detector + tres clasificadores: vertical, horizontal, \textit{quad}).
    \item Todas las reglas de \textit{tracking} temporal (incluyendo la regla de seguridad GREEN→YELLOW→RED, la histéresis para \textit{black} y la detección de \textit{blink} en verde).
    \item Todos los parámetros numéricos extraídos del código (sigmas, pesos, \textit{thresholds}, ventanas temporales).
\end{itemize}

\textbf{Excluir} (infraestructura no funcional):
\begin{itemize}
    \item CyberRT como \textit{framework} de mensajería.
    \item El sistema de \textit{build} Bazel.
    \item La contenerización Docker y los requisitos de hardware específicos.
    \item La integración con otros módulos de Apollo (planificación, control).
\end{itemize}

\textbf{Reemplazar por una interfaz simplificada} (sin afectar la lógica \textit{core}):
\begin{itemize}
    \item El HD-Map dinámico y la proyección 3D→2D, reemplazados por \textit{projection boxes} pre-calculadas con \textit{signal\_id} estable.
    \item El sistema multi-cámara, reducido a una sola cámara.
\end{itemize}

\subsection{Decisiones de Interfaz}

Comparando \textit{inputs}:

\begin{table}[H]
\centering
\begin{tabular}{|l|p{5cm}|p{5cm}|}
\hline
& \textbf{Apollo} & \textbf{Módulo aislado} \\
\hline
Imagen & vía CyberRT (\textit{channel}) & archivo o \textit{frame} de video \\
\hline
\textit{Projection boxes} & calculadas en tiempo real desde el HD-Map y la pose & pre-calculadas en archivo de texto, con \textit{signal\_id} \\
\hline
\end{tabular}
\caption{Comparación de las interfaces de entrada del módulo en Apollo y en el módulo aislado.}
\end{table}

\textbf{Justificación:} desde el punto de vista del resto del módulo TLR, ambas interfaces producen el mismo \textit{input}: una lista de \textit{bounding boxes} 2D con identidad persistente sobre la cual buscar semáforos. La diferencia está en \textit{cómo} se obtiene esa lista, no en \textit{qué representa}. La simplificación no afecta la lógica de detección, asignación, reconocimiento ni \textit{tracking}: estas etapas reciben exactamente la misma estructura de datos.

La salida se mantiene equivalente: un conjunto de semáforos identificados por \textit{signal\_id}, cada uno con su color y \textit{flag} de \textit{blink}.

\subsection{Criterio de Equivalencia Funcional}

Definimos que el módulo aislado es ``equivalente'' al original si:

\begin{enumerate}
    \item \textbf{Mismos parámetros:} todos los valores numéricos (\textit{thresholds}, pesos, sigmas, ventanas temporales) coinciden con los extraídos del código de Apollo.
    \item \textbf{Misma lógica:} cada etapa implementa la misma secuencia de operaciones en el mismo orden, incluyendo casos especiales y validaciones.
    \item \textbf{Mismas reglas:} las reglas de \textit{tracking} y seguridad se aplican idénticamente (GREEN→YELLOW→RED se rechaza; \textit{black} se descarta cuando había color previo; \textit{blink} solo se detecta en verde con el \textit{threshold} de 0.55\,s; el historial se resetea pasados 1.5\,s).
    \item \textbf{Mismos modelos:} se utilizan los mismos pesos de los cuatro modelos pre-entrenados, convertidos al \textit{runtime} de la implementación aislada (PyTorch en lugar de Caffe).
\end{enumerate}

\textbf{Nota importante:} no podemos validar empíricamente la equivalencia ejecutando ambos sistemas en paralelo sobre las mismas entradas: el costo de configurar Apollo es desproporcionado al objetivo (ver Capítulo~\ref{cap:reimplementacion}). La equivalencia se garantiza ``por construcción'': el módulo aislado se deriva directamente del código fuente del original, etapa por etapa, parámetro por parámetro. Esta es una limitación explícita del enfoque, y el conjunto de supuestos bajo los cuales aceptamos la equivalencia se enumera al cierre del Capítulo~\ref{cap:reimplementacion}.

\section{Resumen}

El análisis del módulo TLR estuvo guiado por cuatro preguntas dirigidas al aislamiento. A partir del código fuente se reconstruyó la estructura del módulo (cinco etapas dentro de cuatro subcomponentes en \texttt{modules/perception/}), su interfaz con el resto de Apollo (HD-Map, CyberRT, calibración de cámara, modelos en Caffe), y su comportamiento a alto nivel (\textit{frame} + información geográfica $\to$ semáforos identificados con color y \textit{blink}).

Cuatro hallazgos resultaron determinantes para definir el aislamiento: los dos tipos de identificadores y la importancia del \textit{signal\_id} para el \textit{tracking}, el \textit{semantic\_id} diseñado pero no implementado, la dependencia crítica del HD-Map como fuente de identidad, y la selección automática de cámara. De ellos se desprenden tres decisiones: qué se incluye en el módulo aislado, qué se descarta, y qué se reemplaza por interfaces simplificadas sin alterar la lógica \textit{core}.

El siguiente capítulo entra en el detalle técnico de cada una de las cinco etapas: parámetros exactos, sub-algoritmos (NMS, matriz de costos del húngaro, \textit{prob2color}, reglas de transición), y referencias al código fuente. El Capítulo~\ref{cap:reimplementacion} describe luego cómo se construyó el módulo aislado a partir de estas decisiones, qué cambió funcionalmente respecto al original, y bajo qué supuestos se considera equivalente.